\documentclass[12pt,a4paper]{article}
\usepackage[margin=25mm, headheight=15pt, headsep=10pt]{geometry}
\usepackage{fontspec}
\usepackage[table]{xcolor} % Note: table option is required for rowcolors
\usepackage{titlesec}
\usepackage{tocloft}
\usepackage{fancyhdr}
\usepackage{booktabs}
\usepackage{tcolorbox}
\tcbuselibrary{skins, breakable}
\usepackage[colorlinks=true, linkcolor=emerald, urlcolor=emerald, bookmarks=true]{hyperref}

\definecolor{emerald}{RGB}{0,102,51}
\definecolor{darkred}{RGB}{139,0,0}
\definecolor{lightgray}{RGB}{245,245,245}

\usepackage[nil]{babel}
\babelprovide[import, main]{bengali}
\babelprovide[import]{arabic}
\babelfont{rm}[Renderer=HarfBuzz,Scale=1.0]{Noto Serif Bengali}
\babelfont{sf}[Renderer=HarfBuzz]{Noto Sans Bengali}
\babelfont{tt}[Renderer=HarfBuzz]{Noto Sans Bengali}
\babelfont[arabic]{rm}[Renderer=HarfBuzz,Scale=1.1]{Noto Naskh Arabic}

% Section styling
\titleformat{\section}{\Large\bfseries\color{emerald}}{\thesection.}{0.5em}{}
\titleformat{\subsection}{\large\bfseries\color{emerald!85!black}}{\thesubsection.}{0.5em}{}
\titleformat{\subsubsection}{\normalsize\bfseries\color{emerald!70!black}}{\thesubsubsection.}{0.5em}{}
\titlespacing*{\section}{0pt}{18pt}{8pt}

% Fancy Header/Footer setup
\pagestyle{fancy}
\fancyhf{} % clear all header and footer fields
\fancyhead[L]{\color{emerald}\small\textbf{\nouppercase{\leftmark}}}
\fancyfoot[L]{\scriptsize\color{black!50}{\lat{AI}}-সংকলিত, আলেম-যাচাইকৃত নয়}
\fancyfoot[C]{\color{emerald!70!black}\thepage}
\renewcommand{\headrulewidth}{0.4pt}
\renewcommand{\headrule}{\hbox to\headwidth{\color{emerald}\leaders\hrule height \headrulewidth\hfill}}
\renewcommand{\footrulewidth}{0pt}

% Custom boxes
% 1. Ayat/Hadith Box
\newtcolorbox{ayatbox}[1][]{
    enhanced, breakable, 
    colback=emerald!5, 
    colframe=emerald, 
    boxrule=0.5pt, 
    arc=3pt, 
    left=10pt, right=10pt, top=10pt, bottom=10pt,
    #1
}

% 2. Quote/Translation Box
\newtcolorbox{quotebox}[1][]{
    enhanced, breakable,
    frame hidden,
    colback=lightgray,
    borderline west={3pt}{0pt}{emerald},
    left=15pt, right=10pt, top=10pt, bottom=10pt,
    sharp corners,
    #1
}

% 3. AI-generated notice box
\newtcolorbox{warnbox}[1][]{
    enhanced, breakable,
    colback=red!4,
    colframe=darkred,
    boxrule=0.8pt,
    arc=3pt,
    left=10pt, right=10pt, top=10pt, bottom=10pt,
    #1
}

% Commands
\newcommand{\arab}[1]{\foreignlanguage{arabic}{#1}}
\newcommand{\kib}[1]{\textcolor{emerald}{\textbf{#1}}}

% Latin (English) fallback: Noto Serif Bengali has NO Latin glyphs
\newfontfamily\latfont[Renderer=HarfBuzz]{Noto Serif}
\newcommand{\lat}[1]{{\latfont #1}}

\setlength{\parskip}{5pt}
\setlength{\parindent}{0pt}

% Fix for missing bullet in Noto Serif Bengali
\renewcommand{\labelitemi}{$\bullet$}



\begin{document}
\begin{center}{\Large\bfseries\color{emerald} \lat{05}-সূরা-মিলানো-কিরাআত}\end{center}
\vspace{1em}
\section*{০৫ — সূরা মিলানো (কিরাআত): ফাতিহার পর কী পড়বেন}

\begin{warnbox}
 \textbf{গুরুত্বপূর্ণ নোটিশ:} এই কন্টেন্টটি \lat{AI} দিয়ে প্রস্তুত। কেবল সহীহ/হাসান হাদিস রাখা হয়েছে। আলেম-যাচাইকৃত নয়।
\end{warnbox}

\medskip\hrule\medskip

\subsection*{১. সূত্র কার্ড}

\begin{table}[h]\centering\small
\begin{tabular}{ll}
বিষয় & বিবরণ \\
\hline
\textbf{বিষয়} & সূরা ফাতিহার পর সূরা পাঠ (কিরাআত) \\
\textbf{সূত্র} & সহীহ বুখারি ৭৬২, সহীহ মুসলিম ৪৫১ \\
\textbf{মর্যাদা} & \textbf{সুন্নাতে মুআক্কাদাহ} — প্রথম দুই রাকাতে ফাতিহার পর সূরা পড়া \\
\hline
\end{tabular}
\end{table}

\medskip\hrule\medskip

\subsection*{২. কোন সূরা পড়বেন?}

\subsubsection*{\textbf{প্রথম দুই রাকাতে (ফরজ/সুন্নাত/নফল):}}
\begin{itemize}
\item ফাতিহার পর \textbf{একটি সূরা} পড়ুন।
\item ছোট সূরা (কাওসার, আসর, ইখলাস, ফালাক, নাস, মাউন, কুরাইশ) বা দীর্ঘ সূরার অংশ।
\item প্রতিটি রাকাতে \textbf{ভিন্ন সূরা} পড়া ভালো (ইমাম হলে); একাই নামাজ পড়লে প্রথম দুই রাকাতে ভিন্ন, পরের রাকাতে যা মনে হয়।
\end{itemize}
\subsubsection*{\textbf{তৃতীয় ও চতুর্থ রাকাতে (শাফি'ি):}}
\begin{itemize}
\item শুধু \textbf{সূরা ফাতিহা} — অন্য কিছু নয়।
\end{itemize}
\subsubsection*{\textbf{তাহাজ্জুদ, বিতর, ইশরাক ইত্যাদি নফলে:}}
\begin{itemize}
\item ফাতিহার পর সূরা পড়া \textbf{আবশ্যিক}।
\end{itemize}
\medskip\hrule\medskip

\subsection*{৩. ছোট সূরা তালিকা (নামাজে পঠিত)}

\begin{table}[h]\centering\small
\begin{tabular}{llll}
\# & সূরা & আয়াত সংখ্যা & সংক্ষিপ্ত বিষয় \\
\hline
১ & আল-কাওসার (১০৮) & ৩ & অপার দান/নেয়ামত \\
২ & আল-আসর (১০৩) & ৩ & সময়ের শপথ — মানুষ ক্ষতিতে \\
৩ & আল-ইখলাস (১১২) & ৪ & তাওহীদ — আল্লাহ একক \\
৪ & আল-ফালাক (১১৩) & ৫ & অনিষ্ট থেকে সুরক্ষা \\
৫ & আন-নাস (১১৪) & ৬ & কুমন্ত্রণা থেকে সুরক্ষা \\
৬ & আল-কুরাইশ (১০৬) & ৪ & কুরাইশদের স্থিরতা \\
৭ & আল-মাউন (১০৭) & ৭ & এতীম ও গরীবের প্রতি অবহেলা \\
\hline
\end{tabular}
\end{table}

\medskip\hrule\medskip

\subsection*{৪. কীভাবে মিলিয়ে পড়বেন}

\begin{enumerate}
\item \textbf{ফাতিহা শেষে} কিছুক্ষণ \textbf{বিরতি} দিন।
\item \textbf{"আমিন"} বলুন (নীরবে)।
\item \textbf{"বিসমিল্লাহির রাহমানির রাহীম"} বলুন (আলুদা নয়)।
\item সূরার \textbf{১-৩ আয়াত} পড়ুন।
\item \textbf{"সামিআল্লাহু লিমান হামিদাহ"} বলে রুকুতে যান।
\end{enumerate}
\medskip\hrule\medskip

\subsection*{৫. খুশুর জন্য মূল পয়েন্টস}

\begin{table}[h]\centering\small
\begin{tabular}{ll}
পয়েন্ট & কীভাবে প্রয়োগ করবেন \\
\hline
\textbf{ভিন্নতা} & প্রতিটি রাকাতে ভিন্ন সূরা পড়ার চেষ্টা করুন — এতে মনোযোগ বাড়ে \\
\textbf{অর্থ} & শুধু উচ্চারণ নয় — সূরার অর্থ মনে করুন \\
\textbf{দীর্ঘতা} & নফলে দীর্ঘ সূরা পড়া ভালো; ফরজে সংক্ষিপ্ত রাখুন (ইমাম হলে) \\
\textbf{তাড়াহুড়ো নয়} & প্রতিটি আয়াতের পর ২-৩ সেকেন্ড বিরতি দিন \\
\hline
\end{tabular}
\end{table}

\medskip\hrule\medskip

\subsection*{৬. সংশ্লিষ্ট হাদিস}

\begin{table}[h]\centering\small
\begin{tabular}{ll}
হাদিস & গ্রন্থ \\
\hline
\textbf{\lat{“}ইমাম যা পড়েন, তোমাদেরও তাই পড়ো।\lat{”}} & সহীহ মুসলিম ৪৫১ \\
\textbf{\lat{“}প্রথম দুই রাকাতে সূরা ফাতিহার পর সূরা পড়ো।\lat{”}} & সুনানে আবু দাউদ ১৪৫৮ \\
\textbf{\lat{“}তোমাদের প্রত্যেকে নিজের ইমামের কথা মেনে চলুক।\lat{”}} & সহীহ বুখারি ৭৩৫ \\
\hline
\end{tabular}
\end{table}

\medskip\hrule\medskip

\textit{শেষ আপডেট: আগস্ট ২০২৬}


\end{document}
